\documentclass[11pt]{article}
\usepackage[margin=1in]{geometry}
\usepackage{amsmath,amssymb}
\usepackage[hidelinks]{hyperref}

\title{Product cutoff with simple and repeated single-chain gaps\\[2pt]
  \large Disproof of Conjecture 00000003407}
\author{}
\date{}

\begin{document}
\maketitle

\noindent
The product of independent uniform refresh chains has total-variation
cutoff both when the single-chain gap is simple and when it is repeated.
The two examples use state spaces of sizes two and three, respectively.

\section*{Single-chain spectrum and the product law}

Let $r\in\{2,3\}$, let $\pi$ be uniform on $r$ states, and let $\Pi$ be
the matrix with every row equal to $\pi$. Take the continuous-time
generator $Q=\Pi-I$. The projection $\Pi$ fixes constants and annihilates
the subspace of functions of mean zero. Thus the eigenvalues of $-Q$ are
$0$ with multiplicity one and $1$ with multiplicity $r-1$. In particular,
the gap is simple for $r=2$ and not simple for $r=3$.

Since $\Pi^2=\Pi$, the transition matrix at time $t$ is
$\Pi+e^{-t}(I-\Pi)$. Starting from state $x$, the law is
$\nu_{t,x}=\rho\delta_x+(1-\rho)\pi$, with $\rho=e^{-t}$. The product of
$n$ copies, with generator equal to the sum of the coordinate generators,
has independent coordinates with these laws. Its stationary measure is
$\pi^{\otimes n}$. Relabeling states in each coordinate carries any
starting vector to any other, so all deterministic starts have the same
distance $d_n(t)$ from stationarity in total variation.

\section*{Upper and lower bounds}

For a single coordinate, the likelihood ratio $\nu_{t,x}/\pi$ equals
$1+(r-1)\rho$ at $x$ and $1-\rho$ elsewhere. Its mean square under $\pi$
is $1+(r-1)\rho^2$. Independence multiplies mean squares of the product
likelihood ratio. Thus Cauchy--Schwarz gives
\[
 d_n(t)\le \frac12\sqrt{(1+(r-1)\rho^2)^n-1}
 \le \frac12\sqrt{\exp\bigl(n(r-1)\rho^2\bigr)-1}.
\]

For the lower bound, let $S$ count coordinates equal to their own starting
states. Its stationary mean is $n/r$, and its mean at time $t$ is
$n/r+\Delta$, where $\Delta=n(1-1/r)\rho$. The count has variance at most
$n/4$ under either law, since it is a sum of independent indicators. For
the event $S\ge n/r+\Delta/2$, the stationary probability and the
complementary probability at time $t$ are both at most $n/\Delta^2$ by
Chebyshev's inequality. Hence
\[
 d_n(t)\ge 1-\frac{2}{n(1-1/r)^2\rho^2}.
\]

\section*{Cutoff}

Set $t_n=\tfrac12\log n$. For fixed $c>0$ and sufficiently large $n$, the
two bounds give
\[
 d_n(t_n+c)\le \tfrac12\sqrt{\exp((r-1)e^{-2c})-1},
 \qquad
 d_n(t_n-c)\ge 1-\frac{2e^{-2c}}{(1-1/r)^2}.
\]
The first bound tends to zero as $c\to\infty$, and the second tends to
one. Therefore, for each fixed $0<\varepsilon<1$, the mixing time
at level $\varepsilon$ is $t_n+O_\varepsilon(1)$. Since $t_n\to\infty$,
this establishes cutoff with a bounded window in both examples
in the usual total-variation sense~\cite[Definition~2.1]{ck}. If the
coordinate generators are averaged instead of summed, the times are
multiplied by $n$ and cutoff still holds.

The same estimates apply in discrete time to the simultaneous product
of kernels $aI+(1-a)\Pi$ for any fixed $0<a<1$. In that case $\rho=a^k$
after $k$ steps, the cutoff center is $\log n/(-2\log a)$ up to rounding,
and the single-chain gap $1-a$ again has multiplicity $r-1$.

\section*{Lean verification}

The Lean 4 project verifies the factorizations of the characteristic
polynomials of $r(I-\Pi)$. It computes the squared weights at rational
parameters, proves the product moment identity by enumerating product
weights, and checks the mean separation and Bernoulli variance bound.
The probability inequalities and the limiting argument are given above.
Run \texttt{lake build} in the \texttt{lean} directory to check the
project.

\begin{thebibliography}{1}
\bibitem{ck}
G.-Y.~Chen and T.~Kumagai, \emph{Cutoffs for product chains}.
\href{https://arxiv.org/abs/1701.06665}{arXiv:1701.06665}.
\end{thebibliography}

\end{document}
